% STATUS: written 2026-11-05 (Lee 2015 read in full; excerpts in excerpts/.
% Beidar--Bresar--Chebotar--Makar-Limanov are named through Lee 2015.)
\section{Applications to derivation problems}\label{sec:functional-identities:applications-to-derivations}

The general theory of
Sections~\ref{sec:functional-identities:fi-theory}
and~\ref{sec:functional-identities:d-free} exists largely to settle
derivation questions, and its first clean application is to Jordan
$\sigma$-derivations. Lee used the functional identity method to prove
a two-variable identity that carries the weight of every later argument in
his paper \cite[Section~2, p.~221]{LeeTK2015JordanFI}.

\begin{theorem}[{\cite[Theorem~2.1, p.~221]{LeeTK2015JordanFI}}]\label{thm:fi:lee21}
Let $R$ be a non-commutative prime ring and let $f\colon R\times R\to
Q_{mr}(R)$ be a bi-additive map. Suppose that $f(x,y)[x,y]=0$ for all
$x,y$ in $R$. Then $f=0$.
\end{theorem}

Here $Q_{mr}(R)$ is the maximal right ring of quotients of $R$
\cite[p.~221]{LeeTK2015JordanFI}. Theorem~\ref{thm:fi:lee21} turns any
bi-additive expression annihilated by the commutator into zero
\cite[p.~221]{LeeTK2015JordanFI}. This is exactly the step needed to
promote Jordan-type laws to derivation laws
\cite[p.~221]{LeeTK2015JordanFI}.

\begin{theorem}[{\cite[Theorem~2.2, p.~221]{LeeTK2015JordanFI}}]\label{thm:fi:lee22}
Let $R$ be a non-commutative prime ring with an endomorphism $\sigma$.
Then every Jordan $\sigma$-derivation $\delta$ of $R$ satisfying
$\delta(xyx)=\delta(x)yx+\sigma(x)\delta(y)x+\sigma(xy)\delta(x)$ is a
$\sigma$-derivation.
\end{theorem}

The proof uses Theorem~\ref{thm:fi:lee21}: after linearization the defect
$\delta(xy)-\delta(x)y-\sigma(x)\delta(y)$ is a bi-additive map whose
product with $[x,y]$ vanishes, so it is zero
\cite[p.~221]{LeeTK2015JordanFI}. When the characteristic is not $2$ the
extra identity holds automatically, so every Jordan $\sigma$-derivation of
a noncommutative prime ring with $\operatorname{char}R\ne 2$ is a
$\sigma$-derivation \cite[Corollary~2.3, p.~221]{LeeTK2015JordanFI}. This
generalizes the earlier characterization of Jordan derivations to the
$\sigma$-derivation context \cite[p.~221]{LeeTK2015JordanFI}.

When $\sigma$ is not injective the conclusion strengthens: every Jordan
$\sigma$-derivation is X-inner, meaning $\delta(x)=ax-\sigma(x)a$ for a
fixed $a$ in $Q_{mr}(R)$ \cite[Theorem~2.4, p.~221]{LeeTK2015JordanFI}.
A Jordan-type law on a prime ring becomes an ordinary derivation law once
it is passed through the functional identity Theorem~\ref{thm:fi:lee21}
\cite[p.~221]{LeeTK2015JordanFI}. The $d$-free method of
Section~\ref{sec:functional-identities:d-free} is what lets this happen on
rings where no bare commuting-map argument reaches
\cite[p.~x]{BresarMartindaleChebotar2007FI}.

Vukman and Kosi-Ulbl prove that an additive map $D$ of a suitably torsion free semiprime ring with identity satisfying $2D(x^{m+n+1})=(m+n+1)(x^mD(x)x^n+x^nD(x)x^m)$ for all $x$ is a derivation mapping the ring into its centre, and that $D=0$ when the ring is a noncommutative prime ring~\cite{WlOnSomeEquationsRelatedToDeriva2006}.

For a suitably torsion free semiprime ring $R$ with derivations $D_1,\dots,D_{n+1}$ satisfying $D_1(x^m)x^{k_1+\cdots+k_n}+x^{k_1}D_2(x^m)x^{k_2+\cdots+k_n}+\cdots+x^{k_1+\cdots+k_n}D_{n+1}(x^m)=0$ for all $x$, Fo\v{s}ner, Fo\v{s}ner and Vukman prove that the $D_i$ sum to zero and that the weighted sum $k_1D_2+(k_1+k_2)D_3+\cdots+(k_1+\cdots+k_n)D_{n+1}$ maps $R$ into its centre, and obtain a range-inclusion result for continuous derivations on Banach algebras~\cite{WlAnIdentityWithDerivationsOnRin2008}.

Demir, De Filippis and Arga\c{c} prove for a prime ring $R$ of characteristic other than two with extended centroid $C$ and a noncentral Lie ideal $L$ that two nonzero generalized derivations $F$ and $G$ of $R$ with $[F(u),G(u)]=0$ for all $u$ in $L$ satisfy $G=\lambda F$ for some $\lambda$ in $C$, except possibly when $R$ satisfies the standard identity of degree four~\cite{WlAQuadraticGeneralizedDifferent2020}.

Ebanks gives a unifying framework for equations of the form $\sum_{k=1}^{n}x^{p_k}f_k(x^{q_k})=0$ for additive maps $f_k$ and integers $p_k,q_k$, solves many equations of this type, and states open problems~\cite{WlCharacterizingRingDerivationsO2014}.
